\chapter*{Abstract}

This thesis analyzes the evolution of video game industry 
revenues over the period 2014–2024, with the objective of 
identifying the main macroeconomic, social, and pandemic-related
determinants of sales across a sample of 39 countries.\\
The analysis combines an exploratory descriptive investigation 
with a Generalized Additive Model with random intercepts (GAM + re), which accounts 
for regional heterogeneity across countries and allows for the 
estimation of non-linear relationships between the predictors 
and the response variable.\\
The model includes GNI per capita, internet penetration, 
unemployment rate, and a tensor product smooth interaction 
between year and depression rates, alongside a continent-level 
random effect. The results indicate that purchasing power follows 
a linear trend, while internet penetration, unemployment, and the 
temporal evolution of depression rates operate through complex 
non-linear paths. Geographic heterogeneity across continents 
accounts for a substantial share of the variance in sales.\\
A notable limitation concerns the Stringency Index, which proved to be
non-significant in the final model despite descriptive evidence 
suggesting a short-term spike in sales during the initial COVID-19 
lockdowns. This discrepancy is attributed to the temporal 
aggregation of the index across a three-year period, which likely 
smoothed out the short-term behavioral response.\\
Overall, the findings suggest that structural and psychological 
factors are more persistent determinants of video game demand than 
short-term pandemic shocks, contributing to a broader macroeconomic 
understanding of the global entertainment industry.
This study contributes to the existing literature by investigating
some of the plausible drivers of video game sales,
taking into account various potential factors and analyzing
their statistical significance through a mixed-effects model.
